\section{Quantum storage on every sufficiently large ring}\label{q1:storage}
This section proves the storage half of Theorem~\ref{q1:main}. A recursion over the levels of the hierarchy bounds the probability that a test at the top level fails, and the decoder recovers the qubit at the end of one tick of the top level that fits in the ring; that tick lasts order \(L\) layers, which gives the linear horizon (Lemma~\ref{q1:l18}).

\paragraph{A stored qubit on each nonhalting window.}

The circuit combines the R1 schedule, an R2 selector and the quantum step. One macro round is one native tick in each of the two classical rules and one quantum tick. At the strobe the circuit uses the centrally normalized bank and reads the current native R2 core's chosen data symbol, with \(D\) the absent symbol, at every physical site \(x\): read that site's own native data field. In the proof its residue is \(x\) modulo the fixed physical stride \(D_{\mathrm{stride}}\) of the native R2 field (Lemma~\ref{q1:l11}) and its logical position is the quotient after removing this residue. Multiple strides or offsets are finite constants. The R1 age and address give the local quantum opcode as in Section~\ref{q1:compiler}. A correct selected epoch has a live selector at every required spacetime site. At the end of each macro round the circuit attempts a full local reset of every retained quantum-capable slot when the selected core symbol is \(D\). A wrong dispatch, a reaction to a malformed intermediate register, or an imperfect terminal cut is counted as a direct physical fault of that strobe if the stipulated semantic controls are correct. A complete local primitive-stencil mark covers any flagged wire in the relevant read, dispatch and erase stencil (probability at most \(C\eps\)); it includes hardware, port moves, waits and newly produced scratch up to its discard. A good core normalization on an \(h\)-close bank is read by majority, not by inspecting an arbitrary single redundant copy. For a completely arbitrary strobe state the semantic word is the total decode after normalization, as defined in Section~\ref{q1:selector}.

\textbf{Base noise cover.} Actual effects different from the ideal clean quantum step, including unintended cuts, are contained in a union of marks in integer (physical site, macro time) spacetime. A mark centered at \(z\) of type \(a\) has footprint and primitive-key support in its ball of radius \(r_{a}\ge 1\) and probability at most \(z_{a}\). All marks use one common independent-key space, and nearby marks may share arbitrarily many keys. The types comprise direct microscopic faults, failed R1 tests in Section~\ref{q1:clock} and the initialized failed core tests of Section~\ref{q1:selector}. For the last type the analysis uses the nominal stationary negative-time comparison, whose possible failure to match our finite encoder is charged separately below. A compiler anomaly mark always means the explicit flag count in its stencil, or an upper bounding stencil union, never disagreement depending on a particular incoming state. The prefix rate mismatch can likewise be marked by a uniform primitive draw in the deterministic threshold band. Thus these marks are measurable in the indicated independent keys. A rare wrong decoded selector or wrong clock may be charged even when the actual gate happened to be the intended one. Per site and time there are finitely or countably many types. For the exponent \(d=d_{\mathrm{moment}}\) fixed in Section~\ref{q1:parameters} before the helper exponents, and every \(\delta>0\), the selector construction of Lemma~\ref{q1:l12} for that \(d\), which does not depend on \(\delta\), and a sufficiently small selected rate, reached through the final padding \(J_{\mathrm{pad}}\) (Lemmas~\ref{q1:l12} and~\ref{q1:l16}), give \(\sum _{a} (1+r_{a})^{d} \sqrt{z_{a}}<\delta\), uniformly in the chosen window. The compiler and R1 geometry precede \(d\) and this choice of \(\delta\). All statements persist under one elementary rate splice inside the window, by Section~\ref{q1:selector} and the envelope bounds on primitive R1 and hardware flags.

\textbf{Probability of a bad top operation.} In this paragraph and the next, \(P,T,h,A,d,\delta\) stand for \(P_{q},T_{q},h_{\mathrm{cover}},A_{\mathrm{pair}},d_{\mathrm{moment}},\delta_{\mathrm{mark}}\) of Section~\ref{q1:parameters}; \(C\) denotes a constant fixed with the geometry, and the constant multiplying \(\delta\) in the recursion below is \(C_{\mathrm{rec}}\). The proof uses the monotone bad-box tests of Section~\ref{q1:compiler}, with base \(B_{0}\) true at a site and tick whenever covered by a mark. Counting all marks regardless of whether a quantum gate failed there is conservative. If the rank-\((i+1)\) parent test is bad, two bad rank-\(i\) children can be chosen with distance at least \(8P^{i}\) between the complete sets of primitive locations involved in their descendant \emph{geometric} tests. For a fixed child, complete footprints have halfwidths at most \(hP^{i}\) and halfheights at most \(hT^{i}\), up to one child width or height at the grid boundaries. Two fixed children with footprints closer than \(8P^{i}\) must have grid offsets in each coordinate at most \(2h+10\), since \(T\ge P\). More than \((10h+40)^{2}\) distinct children cannot all be close to the chosen child. The complete parent footprint has at most \(C(PT)^{i}\) primitive lattice points, with the extra constant factor \(PT\) and fixed collars absorbed. Let \(A\) be a bound on the number of ordered child pairs, fixed once the geometry is fixed.

Let \(u_{i}\) bound the supremum of bad-test probabilities over placements. Truncate \(B_{0}\) by deleting types with \(r_{a}>P^{i}\) when estimating one such pair. Each surviving child test is determined by keys within \(2P^{i}\) of its complete geometric footprint: the mark center and then its key support each require up to \(P^{i}\). Hence the two indicators are independent; each truncated probability is at most \(u_{i}\) by monotonicity. The only way the truncated parent differs is for a deleted mark to hit some primitive point of its full footprint. At most \(C(1+r_{a})^{2}\) centers of a given type hit one point. Thus

\[
\begin{gathered}
u_{0} \le  C \delta, \\ u_{i+1} \le  A u_{i}^{2} + C \delta [PT / P^{d-2}]^{i}.
\end{gathered}
\]

The choice \(d=d_{\mathrm{moment}}\) in Section~\ref{q1:parameters}
makes the bracket at most \(1/4\) and gives \(d\ge4\) for the separate
halting argument.
Choose \(0<b_{*}\le 1/(4A)\) and \(\delta\) with both \(C_{\mathrm{rec}}\delta\le b_{*}/4\) and \(C_{\mathrm{path}} \delta<1/2\) for the path-counting constant of Section~\ref{q1:forgetting}. Induction gives \(u_{i}\le b_{*}2^{-i}\) (Lemma~\ref{q1:l18}). At most a fixed \(C_{\mathrm{top}}\) top boxes (one top-local tick and its port neighbors) matter for the protected nonport. They need no wraparound or dependence estimate between them. The union bound and deterministic reduction show that at the end of the top tick, with fixed recursive recovery, the output on the good event has exactly its input logical state jointly with the reference. The failing event has probability at most \(C_{\mathrm{top}} b_{*}2^{-k}\) at rank \(k\). The controller and mark choices are independent of the logical input.

\textbf{Finite geometry and all elementary times.} One rank-\(k\) top tick spans \(H=T^{k}\) macro rounds and a code home of length \(P^{k}\). The encoder chooses a complete, aligned home in an integer chart (zero may be its origin). It puts the message into one ordinary nonport slot of its rank-\(k\) cell, all other top slots in fixed states and all private work initialized as specified by the encoding. It recursively encodes the \(S\) slots in each parent into its child blocks at the prescribed coordinates, with fresh auxiliary slots. At date zero the top on-site instruction is idle by the timetable for every location; the top port permutation does not act on the message. The encoder initializes enough complete adjacent rank-\(k\) homes, using fixed top input for those homes, to cover the full past cone from the output block through \(H\) dates. For some computable constant \(C_{\mathrm{cone}}\) the needed interval, including the complete boundary blocks and cone of every physical controller transition used, has length at most \(C_{\mathrm{cone}}(H+P^{k}+1)\) (Lemma~\ref{q1:l16}). The physical circuit is finite-speed. The rank \(k\) is the largest with \(4 C_{\mathrm{cone}}(T^{k}+P^{k}+1)\le L\); a total fixed fallback is used if no \(k\ge 1\) qualifies. This whole interval is mapped injectively into the ring, with arbitrary legal filler outside it. The encoder initializes R1 and core boundary symbols in this interval as in their respective sources, and prepares the finite nominal helper marginal jointly with error less than \(1/128\) (Lemma~\ref{q1:l15}). These occur on separate physical wires. The unrestricted endpoint encoder may do arbitrarily much finite work. An analysis coupling uses one infinite extension of all raw keys and both mathematical past histories for the entire finite quantum computation. Geometric witnesses can extend beyond the injective cone, since the events being estimated on the ring have already been coupled by their actual past dependence. Disjoint unwrapped witnesses are not treated as independent on a small physical ring.

Write \(D_{\mathrm{layer}}\) for the fixed number of elementary layers per macro round. For \(t\) from 0 to \(D_{\mathrm{layer}} H\) the decoder appends the remaining elementary layers on every received register at the chosen nominal rate \(\bar\eps_{m}\), then applies the fixed bottom-up recovery and extracts the designated nonport. This is a computable channel at any \(t\), even with a partial contact, live ports or a pending factory. In the proof the combined process simply uses the true elementary rate through \(t\) and nominal afterwards. All estimates above hold for that splice. Beyond the horizon the decoder is any fixed total channel. Thus \(F\ge 1-1/128-C_{\mathrm{top}} b_{*}2^{-k}\). Choose a fixed effective \(k_{0}\) with the final term at most \(1/32\), and take a computable \(L_{1}(n)\) large enough that \(k\ge k_{0}\) thereafter. Then \(H\to\infty\) with \(L\), and in fact \(H>c_{1}L\) by maximality of \(k\) and \(T\ge P\) (Lemma~\ref{q1:l18}). The initialization and recovery depend on \(n,L\) (and \(t\) for recovery) and are independent of the actual rate.


\begin{lemma}[The recursion, horizon and fidelity constants]\label{q1:l18}

Let \(A_{\mathrm{pair}}\ge 1\), \(C_{\mathrm{rec}}\ge 1\) and \(\delta=\delta_{\mathrm{mark}}\) be the constants, and \(u_{i}\ge 0\) the bounds on bad-test probabilities, in the already established geometric recursion

\[
\begin{gathered}
u_{0} \le  C_{\mathrm{rec}} \delta, \\ u_{i+1} \le  A_{\mathrm{pair}} u_{i}^{2} + C_{\mathrm{rec}} \delta q^{i}, \\ q = P_{q} T_{q} / P_{q}^{d_{\mathrm{moment}}-2} \le  1/4.
\end{gathered}
\]

If \(0<b_{*}\le 1/(4 A_{\mathrm{pair}})\) and \(C_{\mathrm{rec}} \delta\le b_{*}/4\), then \(u_{i}\le b_{*}2^{-i}\) for all \(i\ge 0\). Choose \(k_{0}\ge 1\) with \(C_{\mathrm{top}} b_{*}2^{-k_{0}}\le 1/32\). Let \(C_{\mathrm{cone}}\) be a positive integer bounding the complete initialized cone as in Lemma~\ref{q1:l16}. If \(k\) is largest with \(4 C_{\mathrm{cone}}(T_{q}^{k}+P_{q}^{k}+1)\le L\), \(H=T_{q}^{k}\), and \(2\le P_{q}\le T_{q}\), then

\[
\begin{gathered}
L < 12 C_{\mathrm{cone}} T_{q} H, \\ \text{with } c_{1} = 1/(12 C_{\mathrm{cone}} T_{q}),\quad c_{1} L < H \le  D_{\mathrm{layer}} H.
\end{gathered}
\]

Suppose further that \(U(n)\) does not halt, that every elementary rate lies in \(I_{n}\), and that the encoder of this section is used, so that the decoder fails only on the helper-coupling failure of Lemma~\ref{q1:l15}(b), of probability less than \(1/128\), or when one of the at most \(C_{\mathrm{top}}\) top tests is bad, each with probability at most \(u_{k}\). Then for \(k\ge k_{0}\) the nominal-rate completion decoder has entanglement fidelity at least

\[
1 - 1/128 - 1/32 = 123/128 \ge  15/16 = 120/128
\]

at every integer elementary time \(0\le t\le \lfloor c_{1} L\rfloor\).

\end{lemma}
\begin{proof}
The base is \(u_{0}\le b_{*}/4\le b_{*}\). If \(u_{i}\le b_{*}2^{-i}\), the recurrence and \(q\le 1/4\) give

\begin{align*}
u_{i+1} &\le (A_{\mathrm{pair}} b_{*}^{2} + C_{\mathrm{rec}} \delta)4^{-i}
 \le (b_{*}/4+b_{*}/4)4^{-i}\\
 &= (b_{*}/2)4^{-i}
 \le b_{*}2^{-(i+1)}.
\end{align*}

The last inequality is \(2^{-2i}\le 2^{-i}\), valid for \(i\ge 0\). This proves the induction with slack at rank zero as well.

Maximality of \(k\) implies \(L<4 C_{\mathrm{cone}}(T_{q}^{k+1}+P_{q}^{k+1}+1)\). Both \(P_{q}^{k+1}\) and 1 are at most \(T_{q}^{k+1}=T_{q} H\), so the claimed strict bound follows by adding the three terms. Multiplication by \(c_{1}\) yields \(c_{1} L<H\). \(D_{\mathrm{layer}}\) is a positive integer, the number of physical elementary layers per macro round, so \(H\le D_{\mathrm{layer}} H\). Consequently every stated elementary read time lies within the decoder's completion horizon.

On the successful joint coupling, initialization loses probability less than \(1/128\) and the union of the at most \(C_{\mathrm{top}}\) relevant top bad tests loses at most \(C_{\mathrm{top}} b_{*}2^{-k}\). These events are independent of the logical input; the deterministic recursive recovery is the same map for every input and preserves its joint state with the reference on success. On failure its entanglement fidelity is nonnegative. Therefore its unconditional overlap is at least one minus these two probabilities, yielding \(123/128\) for \(k\ge k_{0}\). The elementary-time suffix uses nominal noise from the received entire register; its analysis is the allowed single rate splice, so the same inequality holds at a time inside any physical operation. A computable sufficient size cutoff is \(4 C_{\mathrm{cone}}(T_{q}^{k_{0}}+P_{q}^{k_{0}}+1)\), increased for the other finite-chart conditions. This makes \(k\ge k_{0}\) for all larger \(L\). 
\end{proof}



\paragraph{Joint bad-box recursion.}
\emph{Proof idea.} A bad parent test contains two bad children whose complete footprints are far apart. After marks of radius larger than the child scale are deleted, the two children depend on disjoint noise keys, so their failures are independent and cost \(u_i^2\); the deleted large marks cost a tail term controlled by the radius moment. The resulting recursion \(u_{i+1}\le A_{\mathrm{pair}}u_i^2+C_{\mathrm{rec}}\delta_{\mathrm{mark}}q^i\) with \(q\le1/4\) closes at \(u_i\le b_*2^{-i}\) (Lemma~\ref{q1:l18}).

Let \(B_{0}\) be coverage by the direct microscopic, clock and decomposed initialized-selector marks. Each mark's determining keys and covered footprint lie in its advertised ball. Lemmas~\ref{q1:l1}--\ref{q1:l3}, \ref{q1:l6} and~\ref{q1:l12} show that outside this set a physical instruction, including its \(\operatorname{LOAD}\) and terminal cut, equals the clean reference. A parent test counts named children as in Lemma \ref{q1:l7}. If it is bad, two bad child tests have complete descendant geometric footprints separated by at least \(8P_{q}^{i}\). Remove all base types with radius greater than \(P_{q}^{i}\). A remaining type that hits one child's footprint has its center within \(P_{q}^{i}\) of that footprint and every determining key at most another \(P_{q}^{i}\) away. Hence the two truncated indicators use disjoint primitive keys and are independent. Each has probability at most \(u_{i}\) by monotonicity. Union over at most \(A_{\mathrm{pair}}\) ordered pairs costs \(A_{\mathrm{pair}} u_{i}^{2}\).

A deleted type must hit one of the at most \(C_{\mathrm{geom}}(P_{q} T_{q})^{i}\) points of the complete parent footprint. At most \((2r+1)^{2}\le 9(1+r)^{2}\) centers of a radius-\(r\) type can hit one fixed point. Since \(z_{a}\le \sqrt{z_{a}}\), the shared moment bound yields

\[
\sum _{r_{a}>P_{q}^{i}} 9(1+r_{a})^{2} z_{a} \le 9 \delta_{\mathrm{mark}} P_{q}^{-(d_{\mathrm{moment}}-2)i}.
\]

Absorb the geometric point count and the factor nine into \(C_{\mathrm{rec}}=\max (C_{\mathrm{base}},C_{\mathrm{tail}})\). The base union gives \(u_{0}\le C_{\mathrm{rec}} \delta_{\mathrm{mark}}\), and the tail gives

\[
u_{i+1}\le A_{\mathrm{pair}} u_{i}^{2}+C_{\mathrm{rec}} \delta_{\mathrm{mark}} \bigl[P_{q} T_{q}/P_{q}^{d_{\mathrm{moment}}-2}\bigr]^{i}.
\]

Lemma \ref{q1:l18} closes this recursion. The controller and mark choices are independent of the encoded logical input. Lemmas~\ref{q1:l4}--\ref{q1:l7} put successful outputs in one common correctable operator span, so the fixed bottom-up recovery preserves the state with its reference. Lemma \ref{q1:l16}'s actual finite cone, not a globally periodic hierarchy, gives the every-length ring coupling. The total encoder and decoder and the elementary-time nominal suffix are given by Lemma \ref{q1:l15}. The fidelity and linear horizon are exactly Lemma \ref{q1:l18}.
